\documentclass{article}

\usepackage{bm}
\usepackage{ctex}
\usepackage{tikz}
\usepackage{float}
\usepackage{xcolor}
\usepackage{setspace}
\usepackage{datetime}
\usepackage{geometry}
\usepackage{graphicx}
\usepackage{enumitem}
\usepackage{listings}
\usepackage{tcolorbox}
\usepackage{nicematrix}
\usepackage{amsmath, amssymb, amsfonts, mathrsfs}
\usepackage[fontsize = 12pt]{fontsize}

\renewcommand{\thesubsection}{(\alph{subsection})}
\newcommand{\diff}[2]{\dfrac{\mathrm{d}#1}{\mathrm{d}#2}}
\newcommand{\diffn}[3]{\dfrac{\mathrm{d}^{#3}#1}{\mathrm{d}#2^{#3}}}
\newcommand{\piff}[2]{\dfrac{\partial{#1}}{\partial{#2}}}
\newcommand{\eval}[2]{\left.#1\right|_{#2}}
\newcommand{\e}{\mathrm{e}}
\renewcommand{\d}{\mathrm{d}}
\newcommand{\barline}[1]{\overline{#1\rule{0pt}{1.9ex}}}

\geometry{
    a4paper,
    left = 2cm,
    right = 2cm,
    top = 2.4cm,
    bottom = 2.4cm
}

\setitemize[1]{
    itemsep = 7pt,
    partopsep = 0pt,
    parsep = \parskip,
    topsep = 5pt
}

\title{\vspace{-3em}应用统计：第四次作业}
\author{Illusionna \quad 2025..........004 \quad 人工智能学院}
\date{\currenttime,\ \today}



\begin{document}

\maketitle\vspace*{-1.5em}



\section{Q-3.5}
\textit{\textcolor{gray}{设总体 $X$ 服从区间 $[0,\ \theta]$ 上的某种分布，其密度函数为：
\[ P(x, \theta) = \begin{cases} 
    \dfrac{1}{\theta}x, & 0 \leqslant x \leqslant \theta
    \\[7pt]
    0, & others
\end{cases} \]
(1) 设 $X_1, X_2, \cdots, X_n$ 是样本，求参数 $\theta$ 的矩估计量；\\[5pt]
(2) 假定测得样本值 $(13, 6, 17, 22, 3, 11)$，求 $\theta$ 的矩估计值。}}

\paragraph*{参数 $\theta$ 的矩估计量}~{}

计算总体一阶原点矩，即数学期望：
\[ E(X) = \int_{-\infty}^{+\infty} x P(x, \theta) \d x = \int_{0}^{\theta} x \cdot \dfrac{1}{\theta}x \d x = \dfrac{1}{\theta} \int_{0}^{\theta} x^2 \d x = \dfrac{1}{\theta} \cdot \left[ \dfrac{1}{3}x^3 \right]_{0}^{\theta} = \dfrac{1}{3}\theta^2 \]

令一阶样本矩 $\barline{X}$ 等于一阶总体矩 $E(X)$：
\[ \barline{X} = \dfrac{1}{3}\theta^2 \]

由题意 $x \in [0,\ \theta]$ 可知 $\theta \geqslant 0$，解得参数 $\theta$ 的矩估计量为：
\[ \hat{\theta} = \sqrt{3\barline{X}} \]

\paragraph*{参数 $\theta$ 的矩估计值}~{}

已知样本值为 $(13, 6, 17, 22, 3, 11)$，计算样本均值 $\barline{x}$：
\[ \barline{x} = \dfrac{1}{n}\sum_{i=1}^{n}x_i = \dfrac{13 + 6 + 17 + 22 + 3 + 11}{6} = \dfrac{72}{6} = 12 \]

将样本均值代入矩估计量公式，得到 $\theta$ 的矩估计值为：
\[ \hat{\theta} = \sqrt{3\barline{x}} = \sqrt{3 \times 12} = \sqrt{36} = 6 \]



\section{Q-3.6}
\textit{\textcolor{gray}{设 $\hat{\theta}$ 是参数 $\theta$ 的无偏估计量，且有 $D(\hat{\theta}) > 0$，请证明 $\hat{\theta}^2$ 不是 $\theta^2$ 的无偏估计量。}}

\paragraph*{证明}~{}

因为 $\hat{\theta}$ 是参数 $\theta$ 的无偏估计量，则由定义：
\[ E(\hat{\theta}) = \theta \]

由方差：
\[ D(X) = E(X^2) - [E(X)]^2 \]

可得：
\[ D(\hat{\theta}) = E(\hat{\theta}^2) - [E(\hat{\theta})]^2 \]

则：
\[ D(\hat{\theta}) = E(\hat{\theta}^2) - \theta^2 \]

因此 $\hat{\theta}^2$ 的数学期望：
\[ E(\hat{\theta}^2) = \theta^2 + D(\hat{\theta}) \]

已知 $D(\hat{\theta}) > 0$ 所以：
\[ E(\hat{\theta}^2) = \theta^2 + D(\hat{\theta}) > \theta^2 \]

即 $E(\hat{\theta}^2) \neq \theta^2$，所以，$\hat{\theta}^2$ 不是 $\theta^2$ 的无偏估计量。



\end{document}